\documentclass[11pt]{article}
\usepackage[margin=1in]{geometry}
\usepackage{amsmath,amssymb}
\usepackage{xcolor}
\usepackage[colorlinks=true,linkcolor=blue,urlcolor=blue]{hyperref}
\setlength{\parskip}{0.5em}
\setlength{\parindent}{0pt}

\begin{document}

\begin{center}
{\Large\bfseries Custodial Symmetry in Randall-Sundrum}\\[2pt]
{\large A pedagogical walkthrough}\\[4pt]
{\small June 2026}
\end{center}

\vspace{0.5em}
Built up from the ground: where custodial symmetry comes from, how it is added to
RS, what it modifies and why, the electroweak-versus-flavor split, and when the
protection is valid. Each piece is tied to what our code actually does.

\section{Why there is a problem at all}

In Randall-Sundrum the SM gauge fields live in the 5D bulk between a UV brane
(Planck scale) and an IR brane (TeV scale). Every 4D gauge boson is the
\emph{zero mode} of a 5D field, with an infinite tower of Kaluza-Klein (KK)
excitations above it at the TeV scale. Electroweak symmetry breaking happens on (or
near) the IR brane, where the Higgs lives. When the Higgs gets its vev it does not
only give mass to the zero-mode $W$ and $Z$, it also \emph{mixes} them with their
heavy KK partners; integrating those out shifts the properties of the light $W$ and
$Z$. Those shifts are what the LEP/SLC precision data measure, encoded in the
oblique parameters $S,T,U$.

The dangerous one is $T$. It measures the breaking of \emph{custodial isospin},
which controls
\[
\rho = \frac{M_W^2}{M_Z^2\cos^2\theta_W},
\]
measured to be $1$ to high precision. In minimal RS the correction to $T$ is not
just nonzero, it is \textbf{enhanced by the warp volume},
\[
\Delta T_{\min}\simeq +\frac{\pi v^2}{2\cos^2\theta_W M_{KK}^2}\,L,
\qquad L\equiv k\pi r_c\approx 35\text{--}37,
\]
where $L$ is the logarithm of the Planck-to-TeV hierarchy, the very thing RS is
built to generate. It is large and cannot be made small, so $T$ forces the first KK
gauge bosons up to roughly $10$~TeV. This is the ``RS $T$ problem'': a little-hierarchy
problem in which precision electroweak data push the new physics back up to
$\sim\!10$~TeV. It happens because the bulk gauge symmetry of minimal RS is just the
SM's $\mathrm{SU(2)}_L\times\mathrm{U(1)}_Y$, and the IR brane where EWSB happens has
no reason to respect custodial isospin, so nothing forbids a large, volume-enhanced
$T$.

\section{Where the cure comes from: custodial symmetry in the SM}

Step back to the ordinary Standard Model. Write the Higgs doublet as a $2\times2$
matrix $\Sigma=(\tilde H, H)$. The Higgs \emph{potential} is invariant under a larger
global symmetry than one might expect: $\mathrm{SU(2)}_L\times\mathrm{SU(2)}_R$,
acting as $\Sigma\to L\,\Sigma\,R^\dagger$. The vev $\Sigma\to v\,\mathbb{1}$ breaks
this to the diagonal subgroup $\mathrm{SU(2)}_{L+R}\equiv\mathrm{SU(2)}_V$, the
\textbf{custodial symmetry}.

That diagonal symmetry is what guarantees $\rho=1$ at tree level in the SM: the three
would-be Goldstones (the longitudinal $W,Z$) form a custodial triplet, locking the
$W$ and $Z$ masses in the symmetric ratio. $T$ stays zero until something breaks
$\mathrm{SU(2)}_V$; in the SM only the hypercharge coupling $g'$ and the top-bottom
Yukawa splitting do, and both are small. The point: \textbf{custodial
$\mathrm{SU(2)}_V$ forbids $T$.} If you could impose it on the new physics, $T$ would
be protected.

\section{How custodial is added to RS}

The idea (Agashe, Delgado, May, Sundrum 2003, ``ADMS'') is to \emph{gauge} this
custodial symmetry in the bulk, enlarging the bulk gauge group:
\[
\mathrm{SU(2)}_L\times\mathrm{U(1)}_Y\ \longrightarrow\
\mathrm{SU(2)}_L\times\mathrm{SU(2)}_R\times\mathrm{U(1)}_X.
\]
Hypercharge is now a combination, $Y=T^3_R+X$, with the extra $\mathrm{U(1)}_X$
essentially $(B-L)/2$. To recover the SM at low energy the extra symmetry is broken
back with \textbf{boundary conditions}:
\begin{itemize}
\item On the \textbf{UV brane}, $\mathrm{SU(2)}_R\times\mathrm{U(1)}_X\to\mathrm{U(1)}_Y$
(Dirichlet BCs for the broken generators remove their zero modes and give the
$W_R,Z'$ KK modes UV-brane masses, so there is no light right-handed $W$ or extra
$Z$).
\item On the \textbf{IR brane}, the Higgs is a \emph{bidoublet} $(2,2)$ of
$\mathrm{SU(2)}_L\times\mathrm{SU(2)}_R$, and its vev breaks
$\mathrm{SU(2)}_L\times\mathrm{SU(2)}_R\to\mathrm{SU(2)}_V$, exactly as in the SM.
\end{itemize}
Now the IR brane, where EWSB happens, \emph{respects the diagonal custodial
$\mathrm{SU(2)}_V$}, so the large volume-enhanced $T$ is forbidden by the bulk
symmetry. The only custodial breaking left is the $\mathrm{SU(2)}_R\to\mathrm{U(1)}_Y$
breaking on the \emph{UV} brane, far from the IR-localized fields, so the zero modes
barely feel it.

\section{What this does to $T$ (and why $S$ is untouched)}

The result is the closed form our code uses, CGHNP Eq.~153:
\[
\Delta T_{\rm cust}\simeq -\frac{\pi v^2}{4\cos^2\theta_W M_{KK}^2}\,\frac1L.
\]
Compared with the minimal case the $+L$ became $-1/L$: the sign flipped and the
volume \emph{enhancement} became a volume \emph{suppression}. The ratio is
\[
\left|\frac{T_{\min}}{T_{\rm cust}}\right| = 2L^2\approx 2450,
\]
so custodial $T$ is $\sim\!2450\times$ smaller. The leading $L$-enhanced piece is
now symmetry-forbidden; what survives is a small residual from the UV-brane breaking
that the IR-localized zero mode only feels through its exponentially small UV tail,
hence the $1/L$. In our run this pulls the oblique floor from $\sim\!16$ to
$\sim\!5.8$~TeV.

The floor does not go to zero because \textbf{$S$ is not protected}. $S$ measures a
different combination (roughly the $W$-$B$ kinetic mixing from the KK tower), and
custodial $\mathrm{SU(2)}_V$ says nothing about it. So once $T$ is killed, $S$ sets
the residual floor ($\sim\!2$--$3$~TeV in the literature, $\sim\!5.8$~TeV in our more
conservative proxy). Custodial fixes $T$; $S$ survives and sets the wall.

\section{The second problem, and the second tool: $P_{LR}$}

Gauging $\mathrm{SU(2)}_R$ fixes the oblique $T$ but not the $Z\to b\bar b$ vertex.
The top is heavy, so the left-handed doublet $(t,b)_L$ is IR-localized
($c_{Q_3}<1/2$); $b_L$ lives in that doublet, so it is IR-localized too, and
IR-localized fermions overlap strongly with the KK gauge bosons and the Higgs. The
$Z b_L\bar b_L$ coupling therefore gets a large correction $\delta g_L^b$, shifting
$R_b,A_b,A_{FB}^b$ and giving a bound independent of $T$.

The fix needs an additional \emph{discrete} symmetry, $P_{LR}$, the exchange
$L\leftrightarrow R$ inside the custodial group (Agashe, Contino, Da~Rold, Pomarol
2006, ``ACDP''). The dangerous correction comes from the broken generators and is
\emph{odd under $P_{LR}$}, so if $b_L$ is a $P_{LR}$ eigenstate the correction must
cancel. $P_{LR}$ exchanges $T^3_L\leftrightarrow T^3_R$, so the protected point is
\[
T^3_L = T^3_R.
\]
One arranges this by embedding $(t,b)_L$ in a \textbf{bidoublet} $(2,2)_{2/3}$ where
$b_L$ sits at $(T^3_L,T^3_R)=(-\tfrac12,-\tfrac12)$; there $P_{LR}$ leaves it
invariant and $\delta g_L^b\to0$ at tree level. The price is physical: a bidoublet
has four components, so besides $(t,b)_L$ it contains exotic partners, including a
charge-$5/3$ quark and an extra charge-$2/3$ quark. These are the \textbf{custodians},
new TeV-scale KK fermions, and they are where the loop corrections below come from.

For the right-handed $b$: a singlet $(1,1)_{-1/3}$ has $T^3_L=T^3_R=0$, trivially
$P_{LR}$-symmetric, so $\delta g_R^b=0$. (Other embeddings, e.g.\ the $(2,4)$-type used
to \emph{fit} $A_{FB}^b$, break this and give a nonzero $\delta g_R^b$. So
$g_R^b=0$ is an embedding choice, not a theorem.)

\section{The decomposition in our code}

Custodial is therefore two protections doing two jobs, which is the
\texttt{custodial\_rs\_su2r} vs \texttt{custodial\_rs\_plr} split in the code:
\begin{itemize}
\item \textbf{$\mathrm{SU(2)}_R$} protects the oblique $T$: floor $16\to5.8$~TeV. It
does nothing to $Z b_L$.
\item \textbf{$P_{LR}$} protects $Z b_L$: $\delta g_L^b\to0$. It does nothing to the
oblique floor.
\end{itemize}
Run side by side, \texttt{custodial\_rs\_su2r} and \texttt{custodial\_rs\_plr} give
the \emph{same} $5.8$~TeV oblique floor but a \emph{different} $\delta g_L^b$ (the
first leaves it at the minimal $2.6\times10^{-3}$, the second zeroes it). Two
symmetries, two observables, independent.

\section{Electroweak versus flavor}

Custodial is an \textbf{electroweak} symmetry: it gauges
$\mathrm{SU(2)}_R\times\mathrm{U(1)}_X$ and a discrete parity inside it, acting on
electroweak quantum numbers. So the things it protects are electroweak: the oblique
parameters and the $Z$ couplings.

The RS \textbf{flavor} problem comes from somewhere else: the \textbf{KK gluon}.
Anarchic 5D Yukawas plus geometric localization generate flavor-violating quark
couplings to the first KK gluon, and KK-gluon exchange produces $\Delta F=2$
operators; the one that dominates $\varepsilon_K$ is the left-right operator $C_4$,
QCD- and chirally enhanced. Crucially, \textbf{the KK gluon is a color octet and
an electroweak singlet.} It is neutral under $\mathrm{SU(2)}_L,\mathrm{SU(2)}_R,
\mathrm{U(1)}_X$, so the custodial group cannot act on it. Custodializing the model
leaves the leading $\varepsilon_K$ amplitude \emph{exactly unchanged}.

This is why the electroweak floors drop to a few TeV but the $\varepsilon_K$ wall
sits at tens of TeV. It is physics, not a code limitation: Blanke et al.\ computed
the entire $\Delta F=2$ set inside the full custodial model and still found
$M_{KK}\gtrsim18$--$30$~TeV, and Bauer~II states $C_4$ ``is a pure QCD effect
$\ldots$ insensitive to the precise embedding of the electroweak gauge symmetry.''
What moves the $\varepsilon_K$ wall is \textbf{flavor alignment} (RH-down $U(3)$
degeneracy, $V_{5KM}$, 5D-MFV), a different axis. Custodial and alignment are
orthogonal cures for orthogonal problems: custodial for electroweak, alignment for
flavor.

\section{Subtleties: when ``nothing'' is not literally nothing}

The same custodial structure that protects the diagonal $Z b_L$ vertex also reshuffles
the flavor-violating $Z$ couplings. Bauer~II Fig.~13 / Eq.~131 is exactly this: the
left-handed $Z$-mediated FCNCs are suppressed by $L$, but the right-handed ones are
\emph{enhanced} by about a factor of $10$. So custodial genuinely changes the
tree-level $Z$ contribution to $\Delta F=2$. For $\varepsilon_K$ this stays subleading
to the KK gluon, so it does not move the wall; where it matters is rare decays like
$K\to\pi\nu\bar\nu$. Hence the careful statement is ``custodial does not move the
$\varepsilon_K$ wall,'' not ``custodial does literally nothing to flavor.''

\section{When is the protection valid, and what are its limits}

\begin{itemize}
\item \textbf{Tree-level, loop-corrected.} $P_{LR}$ zeroes $\delta g_L^b$ at tree
level, but the custodian top-partners in loops (Carena, Pont\'on, Santiago, Wagner,
``CPSW'') bring back a negative $\delta g_L^b$ correlated with $+\Delta T$.
\item \textbf{Embedding-dependent.} The protection holds only for
$q_L\in(2,2)_{2/3}$ and $b_R\in(1,1)_{-1/3}$; other choices lose it (sometimes on
purpose, to fit $A_{FB}^b$).
\item \textbf{Leading-order in $1/L$.} The custodial $T$ (Eq.~153) is the leading
$1/L$ piece; there are subleading corrections.
\item \textbf{It costs new states.} The bidoublets force the custodian fermions,
including the charge-$5/3$ quark, into the TeV-scale spectrum.
\end{itemize}

For \emph{our} implementation specifically: we do not build the full 5D custodial
model, we impose its leading results as a proxy: swap in the custodial $T$
coefficient, zero $\delta g_L^b$ on the protected mask, zero $\delta g_R^b$ for the
singlet, add the CPSW loops, and leave the KK-gluon $\Delta F=2$ untouched. The four
places where that proxy is softer than a real model (the $\kappa_b/L$ residual, the
$\xi/\rho$ loop-mass inputs, the unpinned $g_R$ embedding, the singlet-only $\Delta T$)
are documented in \texttt{docs/CUSTODIAL\_PROVENANCE.md} and are what a full custodial
build (roadmap~D) would replace with derived numbers.

\vspace{0.5em}
\hrulefill

\textbf{In four sentences.} Minimal RS has a volume-enhanced $T$ problem and a
$Z b_L$ problem, both electroweak. Gauging a bulk custodial $\mathrm{SU(2)}_R$ makes
the IR brane respect custodial isospin, killing the enhanced $T$ and leaving the
unprotected $S$ as the residual floor. Adding the discrete $P_{LR}$ with the bidoublet
embedding protects $Z b_L$ on top of that. None of it touches $\varepsilon_K$, set by
the color-octet, electroweak-singlet KK gluon, so flavor needs its own separate cure:
alignment.

\end{document}
